\begin{frame}{Sum hashing}
  La misma idea, ahora con multiconjuntos.
  Los \(r[i]\) son otra vez una tabla aleatoria.
  \[
    \begin{aligned}
      h(\emptyset) &= 0 \\
      h(\{i\}) &= r[i] \\
      h(A + B) &= h(A) + h(B) \\
      h(A \setminus B) &= h(A) - h(B)
    \end{aligned}
  \]
  La resta corresponde a la diferencia cuando \(B\) está contenido en \(A\), contando repeticiones.

  \medskip
  \(h(\{a,a,b\}) = 2\,r[a] + r[b]\).

  \vspace{0.6em}
  {\footnotesize
  Práctica: TC ARG 2024, día 7, problema F.\\
  \url{https://codeforces.com/gym/104686/problem/F}}
\end{frame}
